% Part I > C. The remedy both sides accept. Order (Kyle, 2026-11-11): quotes
% from all three clusters (formalization is universally acknowledged; pushback
% on sufficiency); AI-assisted formalization in practice (pipeline and what is
% being built); it is cheap and it finds errors; verifiers are where AI gains
% follow. talk-v2, 2026-11-11.
\section{Formalization}

% [1.5 min] The community's own words: one or two quotes per cluster, in the
% cluster colours of the Proofs and Prompts slides (cluster.py orders by
% valence: capable/positive = blue, doubted/negative = orange,
% capable/negative = green). All quotes are verbatim from data/posts.jsonl:
% AGMAI (3065, 26 Oct 2026); Etingof & Gerovitch (2114, 21 Oct); Harlev (757,
% 25 Sep); Dolores-Cuenca (1043, 1 Oct); Weinberger (146, 11 Aug, ellipsis
% drops "an idea that has taken millennia to refine,"); Nesterov (830, 29 Sep);
% Kretschmer (2379, 27 Oct); Stiegler (2020, 13 Oct).
\newcommand{\qwho}{\sffamily\upshape\tiny\color{black!55}}
\begin{frame}{Formal verification has advocates in all three groups}
  \centering
  \begin{tikzpicture}[x=1cm, y=1cm,
      qhead/.style={anchor=north west, text width=4.4cm, align=flush left, inner sep=0pt,
                    font=\sffamily\scriptsize\bfseries},
      qbody/.style={anchor=north west, text width=4.4cm, align=flush left, inner sep=0pt,
                    font=\rmfamily\itshape\scriptsize, text=black!85},
      ]
    \useasboundingbox (0,0) rectangle (14.3,-5.9);
    % ---- three clusters ------------------------------------------------------
    \foreach \x/\col/\name in {0/clPos/{AI capable, positive}, 4.85/clNeg/{AI capable, negative}, 9.7/clDoubt/{AI doubted, negative}} {
      \fill[\col] (\x,0) rectangle ++(4.55,-0.07);
      \node[qhead, text=\col] at (\x,-0.2) {\name};
    }
    \node[qbody] at (0,-0.65) {%
      ``Proofs should be Lean-formalized where possible, unless it leads to unacceptable delays.''\par
      {\qwho AGMAI, release guidelines}\par\smallskip
      ``Mathematicians will use Lean formalization and verification of proofs, but will also complement it with human-readable exposition.''\par
      {\qwho Etingof \& Gerovitch, MIT}};
    \node[qbody] at (4.85,-0.65) {%
      ``I expect verifying correctness will soon become a non-issue through a combination of better AI and Lean verification.''\par
      {\qwho Harlev, Cornell}\par\smallskip
      ``\ldots{} an AI that formalizes the proof can further filter the papers.''\par
      {\qwho Dolores-Cuenca, Yonsei, on paper mills}};
    \node[qbody] at (9.7,-0.65) {%
      ``\ldots{} we might be able to solve the problem (that of the convincing incorrect argument) because we have the idea of a proof, \ldots{} and we can insist on formalization.''\par
      {\qwho Weinberger, Chicago}\par\smallskip
      ``\ldots{} many mathematicians would begrudgingly accept its correctness if some AI agent produces suitable Lean code tomorrow.''\par
      {\qwho Nesterov, Stanford, on the finite simple groups}};
    % ---- where they push back ----------------------------------------------
    \draw[black!15, line width=0.4pt] (0,-4.05) -- (14.3,-4.05);
    \node[anchor=north west, inner sep=0pt, font=\sffamily\scriptsize\bfseries, text=black!70] at (0,-4.2)
      {Where they push back};
    \node[qbody, text width=6.6cm] at (0,-4.6) {%
      ``In mathematics, we do judge submissions for non-scientific factors that concern plagiarism, ownership, and originality.''\par
      {\qwho Kretschmer, UT Austin, against judging papers by their Lean certificates alone}};
    \node[qbody, text width=6.6cm] at (7.4,-4.6) {%
      Their initial goal was ``to formalize mathematics in a way that makes it globally understandable.''\par
      {\qwho Stiegler, Paris-Saclay, on friends who work on Lean and are unhappy with where it is heading}};
  \end{tikzpicture}
  \slidesources{Proofs \& Prompts posts 3065, 2114, 757, 1043, 146, 830, 2379, 2020; clusters as on the sentiment slides}
\end{frame}

% [1.5 min] Introduces AI-assisted formalization: the pipeline, then what is
% already being built. The "addresses / stays open" lists that were here moved
% to the scorecard at the end of Part I. Sources: Palomar (Tao, 18 Aug 2026);
% Comparator (Lean FRO, research/02 F); Guo's proposed journal (Tao's blog,
% 17 Oct 2026); conceptual track (Kothari, Meka & Raghavendra, Shtetl-Optimized,
% 19 Sep 2026); OpenAI 42% of 719 (research/03).
\begin{frame}{AI-Assisted Formalization in Practice}
  \centering
  \begin{tikzpicture}[x=1cm, y=1cm, font=\sffamily\footnotesize,
      build/.style={draw=black!15, fill=white, rounded corners=2pt, line width=0.4pt,
                    inner sep=0.17cm, anchor=north west, font=\sffamily\scriptsize,
                    drop shadow={shadow xshift=0.5pt, shadow yshift=-0.7pt, opacity=0.2}}]
    \useasboundingbox (0,0.1) rectangle (14.3,-5.9);
    % ---- definition --------------------------------------------------------
    \node[anchor=north west, text width=14.2cm, align=flush left, inner sep=0pt, font=\sffamily\small]
      at (0,0.05)
      {AI translates mathematics into a proof assistant such as Lean and fills in the proofs.
       A small trusted kernel checks every step. Humans write, or review, what is being proved.};
    % ---- pipeline ----------------------------------------------------------
    \node[stage=PennBlue, minimum width=2.9cm, minimum height=1.05cm] (p1) at (1.65,-1.5)
      {Statement \& blueprint\\[-1pt]{\scriptsize\color{black!60}human}};
    \node[stage=tlPurple, minimum width=2.9cm, minimum height=1.05cm] (p2) at (5.35,-1.5)
      {Formalize \& prove\\[-1pt]{\scriptsize\color{black!60}AI agents}};
    \node[stage=PennGreen, minimum width=2.9cm, minimum height=1.05cm] (p3) at (9.05,-1.5)
      {Kernel check\\[-1pt]{\scriptsize\color{black!60}Lean}};
    \node[stage=PennGold, minimum width=2.9cm, minimum height=1.05cm] (p4) at (12.75,-1.5)
      {Certificate\\[-1pt]{\scriptsize\color{black!60}registry, e.g.\ Palomar}};
    \foreach \u/\v in {p1/p2, p2/p3, p3/p4} \draw[flow] (\u) -- (\v);
    \draw[flow=PennBlue, dashed] (p3.south) .. controls +(0,-0.55) and +(0,-0.55) .. (p1.south);
    \node[note, fill=white, inner sep=1pt] at (5.35,-2.48) {a human checks that the formal statement matches the paper};
    % ---- what is being built ------------------------------------------------
    \node[anchor=north west, inner sep=0pt, font=\sffamily\small\bfseries, text=black!70] at (0,-3.0)
      {Already being built};
    \node[build] at (0.0,-3.45)
      {\parbox[t][1.95cm][t]{3.05cm}{\raggedright {\bfseries\footnotesize Palomar}\par\smallskip
       A registry of Lean-verified mathematics (Lean FRO and ICARM, Aug 2026). Two independent kernels replay every proof.}};
    \node[build] at (3.62,-3.45)
      {\parbox[t][1.95cm][t]{3.05cm}{\raggedright {\bfseries\footnotesize Comparator}\par\smallskip
       Checks that a proof proves exactly the stated theorem, with no extra axioms. Part of Lean since v4.35; used for FLT.}};
    \node[build] at (7.24,-3.45)
      {\parbox[t][1.95cm][t]{3.05cm}{\raggedright {\bfseries\footnotesize New venues}\par\smallskip
       A journal that requires Lean-verified proofs (Guo, proposed Oct 2026). A STOC/FOCS/SODA track where a Lean certificate stands in for the proof (proposed).}};
    \node[build] at (10.86,-3.45)
      {\parbox[t][1.95cm][t]{3.05cm}{\raggedright {\bfseries\footnotesize Lab releases}\par\smallskip
       OpenAI shipped Lean proofs with 42\% of its 719 manuscripts. FLT and the finite simple groups shipped challenge files.}};
  \end{tikzpicture}
  \slidesources{Tao, \emph{Palomar}, blog, 18 Aug 2026 \textbullet{} Lean FRO, Comparator \textbullet{} Guo, guest post on Tao's blog, 17 Oct 2026 \textbullet{} Kothari, Meka \& Raghavendra, Shtetl-Optimized, 19 Sep 2026 \textbullet{} research notes 01--03}
\end{frame}

% [1.5 min] Merges talk-v2's "Formalization got cheap" and "Certificates catch
% what review missed" (full versions in the backup). Left: log scale in days,
% x = 1.05cm * log10(days). Human: Strong PNT 18 months (Tao--Kontorovich);
% FLT 5-year EPSRC grant from 2024 (unfinished); Odd Order theorem alone,
% 6 years (ITP 2013). AI: Strong PNT 3 weeks (Math, Inc.); FLT 11 days
% (Anthropic); CFSG 24 days (Axiom), all self-reported (research/01 s1).
% Cost (Kyle asked for data on cheapness, 2026-11-11): Meta's textbook
% scaffold, arXiv 2604.03071 (Grinberg, Introduction to Algebraic
% Combinatorics, ~130k lines, ~$100k, ~$200/page; research/02 D); FLT ~6B
% output tokens, GBP 91k-226k at list price (AcadMathSci B.1, report
% estimate); fixed capability 5-10x cheaper a year (AcadMathSci B.3, citing
% Gundlach et al., arXiv 2511.23455). Not on the slide: Prove2Me FAQ, a
% $200/month plan covers roughly a textbook; Buzzard, 1.2M lines in three
% weeks against 2.3M for Mathlib in nine years (AcadMathSci A.5).
% Right: OpenAI audit is independent; CFSG and Stacks are Axiom's own reports.
\begin{frame}{Cheap, and it catches errors}
  \centering
  \begin{tikzpicture}[x=1cm, y=1cm, font=\sffamily\scriptsize,
      hbar/.style={fill=black!28, draw=none},
      abar/.style={fill=PennBlue, draw=none},
      blabel/.style={anchor=west, inner sep=0pt, text=black!75},
      rowlab/.style={anchor=south west, inner sep=0pt, font=\sffamily\footnotesize\bfseries},
      colhd/.style={anchor=north west, inner sep=0pt, font=\sffamily\small\bfseries, text=black!70},
      ehead/.style={anchor=north west, text width=6.2cm, align=flush left, inner sep=0pt,
                    font=\sffamily\footnotesize\bfseries},
      ebody/.style={anchor=north west, text width=6.2cm, align=flush left, inner sep=0pt,
                    font=\sffamily\scriptsize, text=black!80}]
    \useasboundingbox (0,0.1) rectangle (14.3,-5.9);
    % ================= left: time and cost to formalize =================
    \node[colhd] at (0,0.05) {Time and cost to formalize};
    \fill[hbar] (0,-0.45) rectangle ++(0.3,-0.18);
    \node[blabel] at (0.38,-0.54) {human-led};
    \fill[abar] (1.85,-0.45) rectangle ++(0.3,-0.18);
    \node[blabel] at (2.23,-0.54) {AI agents (self-reported)};
    \node[note, anchor=west] at (5.75,-0.54) {log scale};
    % rows: y / label / human cm / human text / AI cm / AI text
    \foreach \y/\lab/\h/\htext/\a/\atext in {
        -1.1/{Strong Prime Number Theorem}/2.88/{18 months}/1.39/{3 weeks (Math, Inc.)},
        -2.05/{Fermat's Last Theorem}/3.42/{5-year grant, unfinished}/1.09/{11 days (Claude agents)},
        -3.0/{Finite simple groups}/3.51/{6 years, Odd Order theorem alone}/1.45/{24 days, whole classification (Axiom)}} {
      \node[rowlab] at (0,\y) {\lab};
      \fill[hbar] (0,\y-0.05) rectangle (\h,\y-0.3);
      \node[blabel] at (\h+0.1,\y-0.175) {\htext};
      \fill[abar] (0,\y-0.36) rectangle (\a,\y-0.61);
      \node[blabel] at (\a+0.1,\y-0.485) {\atext};
    }
    % log axis in days
    \draw[black!40, line width=0.5pt] (0,-3.72) -- (3.6,-3.72);
    \foreach \v/\t in {0/1 day, 0.89/1 wk, 1.55/1 mo, 2.69/1 yr, 3.42/5 yr} {
      \draw[black!40] (\v,-3.72) -- (\v,-3.8);
      \node[note, anchor=north] at (\v,-3.82) {\t};
      \draw[black!8, line width=0.4pt] (\v,-3.72) -- (\v,-0.95);
    }
    % cost
    \foreach \x/\big/\cap in {
        0/{\$200 a page}/{Grinberg's 500-page textbook: all 340 targets in about a week, about \$100k (Meta)},
        2.5/{\pounds91--226k}/{Fermat's Last Theorem: about 6B output tokens (AcadMathSci estimate)},
        5.0/{5--10$\times$}/{cheaper each year for a fixed level of capability}} {
      \node[anchor=north west, inner sep=0pt, font=\sffamily\bfseries\large, text=PennBlue] at (\x,-4.38) {\big};
      \node[anchor=north west, inner sep=0pt, text width=2.3cm, align=flush left, font=\sffamily\tiny, text=black!70] at (\x,-4.93) {\cap};
    }
    % ================= right: errors found =================
    \draw[black!12, line width=0.5pt] (7.55,0.05) -- (7.55,-5.6);
    \node[colhd] at (7.9,0.05) {Errors found by formalizing};
    \node[ehead] at (7.9,-0.55) {OpenAI's 719 manuscripts};
    \node[ebody] at (7.9,-0.95)
      {3 withdrawn within a day after a sign error; \textbf{all three were outside the Lean coverage}.
       An independent audit of the 245 formalized families found no theorem shown false.};
    \node[ehead] at (7.9,-2.25) {Finite simple groups};
    \node[ebody] at (7.9,-2.65)
      {34 deviations from the published record, including a false equality in a published
       quasithin lemma (a counterexample in SL$_2(5)$). {\color{black!50}Axiom, self-reported; the build has not been reproduced.}};
    \node[ehead] at (7.9,-3.95) {The Stacks Project};
    \node[ebody] at (7.9,-4.35)
      {Dozens of statements that fail in the stated generality: missing hypotheses, edge cases in
       positive characteristic. {\color{black!50}Axiom, self-reported.}};
  \end{tikzpicture}
  \slidesources{Math, Inc., Strong PNT (2025) \textbullet{} Anthropic, FLT (Sep 2026) \textbullet{} Meta, arXiv 2604.03071 \textbullet{} AcadMathSci, Oct 2026 (B.1, B.3) \textbullet{} Axiom, CFSG (Oct 2026) and Stacks (Nov 2026) \textbullet{} Gonthier et al., ITP 2013 \textbullet{} borism/openai-math-slop, issue \#2 \textbullet{} research note 01}
\end{frame}

% [1.5 min] No verifier, no flywheel (AcadMathSci A.5; the report's view, not
% a measurement). For this audience: AI improves fastest where outputs can be
% checked automatically, because the check is a reward signal. AlphaProof
% figures from A.5. The Yue et al. mechanism (RLVR sharpens rather than
% extends) is still on the verify list, so it is not on the slide.
\begin{frame}{No verifier, no flywheel}
  \centering
  \begin{tikzpicture}[x=1cm, y=1cm, font=\sffamily\footnotesize]
    \useasboundingbox (0,0) rectangle (14.3,-5.9);
    % ---- the loop -----------------------------------------------------------
    \node[stage=tlPurple, minimum width=2.6cm] (m) at (1.75,-1.1) {Model\\[-1pt]{\scriptsize\color{black!60}proposes a proof}};
    \node[stage=black, minimum width=2.6cm] (f) at (5.35,-1.1) {Formal proof\\[-1pt]{\scriptsize\color{black!60}attempt}};
    \node[stage=PennGreen, minimum width=2.6cm] (k) at (3.55,-3.15) {Kernel check\\[-1pt]{\scriptsize\color{black!60}pass or fail}};
    \draw[flow] (m) -- (f);
    \draw[flow] (f) -- (k);
    \draw[flow=PennGreen] (k) -- node[note, left=3pt, align=right] {reward:\\training signal} (m);
    \node[note, anchor=north, text width=6.2cm, align=center] at (3.55,-4.0)
      {AlphaProof autoformalized about 1M problems into about 80M Lean statements to train on this loop.};
    % ---- the argument -------------------------------------------------------
    \node[anchor=north west, text width=6.6cm, align=flush left, inner sep=0pt, font=\sffamily\small] at (7.6,-0.2) {%
      \begin{itemize}\setlength{\itemsep}{5pt}
        \item Where a result can be checked automatically, every attempt is a training signal. AcadMathSci: Lean is becoming the way AI checks its own work.
        \item Where it can't (modelling choices, significance, exposition), a person has to judge every output.
        \item So capability ``will stay uneven across the mathematical sciences''.
      \end{itemize}};
    % key claim (Kyle's wording, 2026-11-11)
    \node[keybox, anchor=north] at (7.15,-4.95) {Where formalization lies, AI-driven gains follow.};
  \end{tikzpicture}
  \slidesources{Reisinger et al., AcadMathSci, Oct 2026 (\S1.2; A.5, the report's assessment)}
\end{frame}
